\begin{answers}{ch:gpu}{なぜ GPU が必要か：計算の正体と高速化}
\item 合計は 1 トークンあたり約 $2.85\times10^{8}$ FLOPs で、ブロックの部分は 72.9\% です。ミニGPT の 71.4\% より少し大きくなります。
\coderef{code:gpu-theory}の式に $V=50257$、$T=1024$、$L=12$、$d=768$ を入れます。
QKV は $L\cdot6d^2\approx4{,}247$ 万（14.9\%）、注意のスコアは $L\cdot4Td\approx3{,}775$ 万（13.3\%）です。
射影は $L\cdot2d^2\approx1{,}416$ 万（5.0\%）、MLP は $L\cdot16d^2\approx1.132$ 億（39.8\%）です。
出力ヘッドは $2dV\approx7{,}719$ 万（27.1\%）です。
ミニGPT では出力ヘッドが 28.6\% だったので、ブロックの割合は 1.5 ポイント増えました。
増えたのは主に注意のスコアです（7.1\% から 13.3\%）。
文脈長 $T$ が 8 倍になったのに、$d$ は 4 倍にしかなっていないからです。
\item 約 10.5 時間です。
総計算量は式\eqref{eq:gpu-6n}より $C\approx6ND=6\times10^{9}\times2\times10^{10}=1.2\times10^{20}$ FLOPs です。
実効の速さは、公称値に枚数と MFU を掛けて $8\times989\times10^{12}\times0.4\approx3.16\times10^{15}$ FLOP/s です。
割ると $1.2\times10^{20}\div3.16\times10^{15}\approx37{,}900$ 秒で、約 10.5 時間になります。
GPU 1 枚あたりに直すと、約 84 GPU 時間です。
\item 長方形の行列積は、1 スレッドで約 53 GFLOP/s でした。大きな正方形（$n=512$ で 62、$n=2048$ で 63）より 15\% ほど遅く、$n=192$ の正方形（51）とほぼ同じです（実測）。
FLOPs は式\eqref{eq:gpu-matmul-flops}の $2mkn=2\times4096\times192\times192\approx3.0\times10^{8}$ で数えました。
2 スレッドでは、長方形が 96、$n=192$ が 109、$n=512$ が 115、$n=2048$ が 127 GFLOP/s でした。
速さは、行列の全体の大きさよりも、短い辺（ここでは 192）で決まっているように見えます。
読み書きするデータ量に対して計算が少ない形ほど、メモリの出し入れが足を引っぱるからだと考えられます。
この形は、ミニGPT の線形層（$32\times128=4096$ トークン、$d=192$）の計算そのものです。
ミニGPT の行列積は、この CPU の最高速の 8 割ほどで動いていることになります。
\item 合計を FP32 で持てば、正しく求まります。
合計だけを FP32 にして BF16 の 0.001 を 1 万回足すと、9.9945 になりました（実測）。
10 との差は、BF16 では 0.001 が 0.00099945 に丸められることによるものです。
合計を BF16 で持つと、合計が大きくなるにつれて隣の数との間隔が広がり、0.001 が丸めで消えてしまいます。
FP32 なら間隔が十分に細かいので消えません。
これが、混合精度学習で重みの更新を FP32 の主重みに対して行う理由です。
BF16 のままでも、同じ大きさの数どうしを 2 つずつ足していけば、結果は 10.0 になります（実測）。
大きな数に小さな数を足さないことが要点です。
\item GPU の計算の完了を、毎ステップ待つことになるからです。
GPU の命令は非同期に実行されます。
Python は計算を頼むだけで先へ進み、GPU が計算している間に次のバッチの準備や次の命令の発行をします。
\code{loss.item()} は、損失の値を CPU に取り寄せます。
値はまだ計算の途中かもしれないので、そこまでの計算がすべて終わるまで Python が止まります。
運ぶデータは数 1 個でも、待ち時間のせいで GPU と CPU が同時に働けなくなるのです。
損失の表示は 100 ステップに 1 回などに減らすと、この待ちをほぼなくせます。
\end{answers}
